\documentclass[12pt, letterpaper]{article}

%-------Packages---------
\usepackage{amssymb,amsfonts}
\usepackage{mathptmx}
\usepackage{amsthm}
\usepackage{graphicx}
\usepackage[all,arc]{xy}
\usepackage[colorlinks=true, allcolors=blue]{hyperref}
\usepackage{enumerate}
\usepackage{bbm}
\usepackage{dsfont}
\usepackage{mathrsfs}
\usepackage{tikz-cd}
\usepackage{setspace}
\usepackage{import}
\usepackage[utf8]{inputenc}
\usepackage{hyperref}
\usepackage{tikz}
\usepackage{float}
\usepackage{mdframed}
\usepackage{fancyhdr}
\usepackage[nointegrals]{wasysym}

\usepackage[left=1in,top=1in,right=1in,bottom=1in]{geometry}

\usepackage{mathtools}
\usepackage{setspace}
\usepackage{cleveref}
\usepackage{multicol}

%--------Theorem Environments--------
%theoremstyle{plain} --- default
\newmdtheoremenv{theorem}{Theorem}[subsection]
\newmdtheoremenv{corollary}[theorem]{Corollary}
\newtheorem{proposition}[theorem]{Proposition}
\newmdtheoremenv{lemma}[theorem]{Lemma}
\newtheorem{conj}[theorem]{Conjecture}
\newtheorem{quest}[theorem]{Question}

\theoremstyle{definition}
\newmdtheoremenv{definition}[theorem]{Definition}
\newmdtheoremenv{definitions}[theorem]{Definitions}
\newtheorem{example}[theorem]{Example}
\newtheorem{algorithm}[theorem]{Algorithm}
\newtheorem{rmk}[theorem]{Remark}
\newtheorem{exercise}[theorem]{Exercise}

%----------- my commands -------------
\newcommand{\C}{\mathbb{C}}
\newcommand{\R}{\mathbb{R}}
\newcommand{\Q}{\mathbb{Q}}
\newcommand{\Z}{\mathbb{Z}}
\newcommand{\N}{\mathbb{N}}
\renewcommand{\P}{\mathbb{P}}
\newcommand{\contr}{\Rightarrow \Leftarrow}
\newcommand{\HRule}[1]{\rule{\linewidth}{#1}}
\newcommand{\s}{\(\,\)}
\renewcommand{\t}[1]{\text{#1}}
\newcommand{\ang}[1]{\left\langle #1 \right\rangle}

% Meta data
\setstretch{1.2}

\begin{document}

\pagestyle{fancy}
\fancyhf{}
\fancyhead[LOE]{\slshape{\textbf{Vincent Ferrara}}}
\fancyhead[ROE]{\thepage}

\subsection*{Problem 1}
\begin{enumerate}[(a)]
    \item Let $\sigma \in S_n$, write $\sigma=a_1b_1\dots a_mb_m$ where $a_i,b_i$ are transpositions\\
    Consider $a_kb_k$\\
    $a_kb_k=(ab)(cd)$ s.t. $a,b,c,d \in \{1,2,\dots, n\}$\\
    $(ab)(cd)=(acb)(acd)$\\
    Thus since we have an even amount of transpositions we can write them all as 3-cycles.
    \item Let $(abc)$ and $(dfg)$ be two 3 cycles in $A_n$\\
    Consider $(abc)(dfg)(abc)^{-1}=(abc)(dfg)(acb)=(dfg)$
    \item \begin{enumerate}[i.]
        \item Let $N \unlhd A_n$ if $N=\{e\}$ then done so assume $N \neq \{e\}$\\
        Let $e \neq \sigma \in N$ s.t. it has the most fixed points.
        \item Let $\sigma=(a_1b_1)(a_2b_2)\dots(a_tb_t)$ bet disjoint transpositions.\\
        $t$ is not $1$ since this would mean $\sigma$ is odd so $t\geq 2$\\
        Since these are distinct this means we fix $n-2t \geq n-4$ elements, and since $n \geq 5 \implies n-2t \geq n-4 \geq 1$ thus there is at least one element that is fixed by $\sigma$\\
        Let $r \in \{1,\dots,n\}$ s.t. $r$ is fixed in $\sigma$\\
        Let $\tau=(a_1b_1r)$\\
        Consider $\tau\sigma\tau^{-1}\sigma^{-1}$\\
        Since $N$ is normal we know $\tau\sigma\tau^{-1} \in N$ and since $\sigma^{-1} \in N$ we know $[\tau, \sigma] \in N$\\
        $=\tau(a_1b_1)\tau^{-1}\tau(a_2b_2)\tau^{-1}\dots\tau(a_tb_t)\tau^{-1}\sigma$\\
        $=(b_1r)(a_2b_2)\dots(a_tb_t)\sigma$\\
        $=(b_1r)(a_1b_1)=(a_1rb_1)$\\
        Thus $[\tau, \sigma]$ fixes $n-3$ points\\
        $n-2t \leq n-4 < n-3$ thus the commutator fixes more points and is not the identity, thus $\sigma$ cannot be written as disjoint transpositions.
        \item Since 3-cycles cant be written as two disjoint transpositions if $\sigma$ is a 3 cycle then we are done so assume not\\
        The only way for a permutation to move an even amount of elements if it can be written as disjoint transpositions thus it $\sigma$ cannot move only 4 elements then it could
        be written as $2$ disjoint transpositions thus $\sigma$ must move at least 5 elements.
        \item Thus $\sigma=(a_1a_2a_3\dots)(\dots)\dots(\dots)$\\
        Let $r,s$ be two elements not in $\{a_1,a_2,a_3\}$ that are moved by $\sigma$.\\ 
        Let $\tau = (a_3rs)$\\
        Consider $\tau\sigma\tau^{-1}\sigma^{-1}$\\
        Since $N$ is normal we know $\tau\sigma\tau^{-1} \in N$ and since $\sigma^{-1} \in N$ we know $[\tau, \sigma] \in N$\\
        Let $x$ be fixed by $\sigma$ then $\tau\sigma\tau^{-1}\sigma^{-1}(x)=\tau\sigma\tau^{-1}(x)$ and since $\tau$ only touches things that are moved by $\sigma$ thus $=\tau\sigma(x)=\tau(x)=x$\\
        Thus if $x$ is fixed by $\sigma$ it is fixed by $[\tau, \sigma]$\\
        Consider $a_2$, $\sigma$ moves $a_2$ now consider $\tau\sigma\tau^{-1}\sigma^{-1}(a_2)=\tau\sigma\tau^{-1}(a_1)=\tau\sigma(a_1)=\tau(a_2)=a_2$\\
        Thus the commutator fixes $a_2$ thus this is a contradiction thus $\sigma$ must be a 3-cycle\\
        Thus for all 3 cycles $x \in A_n$ we have that $\sigma x \sigma^{-1}$ is also a 3 cycle and in $N$ since $N$ is normal.\\
        Thus let $\sigma=(abc)$ and let $x=(dfg)$\\
        $(abc)(dfg)(acb)=(dfg)$ thus we have all 3 cycles in $N$ and since $A_n$ is generated by 3-cycles thus $N=A_n$
    \end{enumerate}
\end{enumerate}

\newpage
\subsection*{Problem 2}
$|S_p|=p!$ thus the order of a $p$-Sylow is just $p$\\
Thus the number of distinct $p$ cycles is $(p-1)!$ thus there are $(p-1)!$ groups generated by each of these $p$ cycles\\
However, proven in last homework since in these groups are of order $p$ we know it has $p-1$ generators thus we are overcouting by a factor of $p-1$ since we count each group $p-1$ times thus $n_p=\dfrac{(p-1)!}{p-1}=(p-2)!$\\
Thus by Sylow 3 we know $n_p=(p-2)! \equiv 1 \pmod{p} \implies (p-1)! \equiv (p-1) \pmod{p} \implies (p-1)! \equiv -1 \pmod{p}$

\newpage
\subsection*{Problem 3}
Consider $SL_2(F_5) = \{\t{2x2 matrices } | \t{ det} = 1\}$\\
Thus for matrix $\begin{pmatrix}
    a & b\\
    c & d
\end{pmatrix}$ in $SL_2(F_5)$ we know $ad-bc=1$\\
Consider the first column we have $5^2-1$ choices since all it has to be is non zero\\
Then for the second column we have $5$ choices for $b$ but then $d$ is fixed to satisfy $ad-bc=1$ thus we have $120$ choices thus $|SL_2(F_5)|=120$\\
Thus since $PSL_2(F_5)=\dfrac{SL_2(F_5)}{\{\pm I\}}$ we know $|PSL_2(F_5)|=\dfrac{|SL_2(F_5)|}{|\{\pm I\}|}=60$\\
Let $x$ and $y$ be two classes of matrices
\[
A=\begin{pmatrix}
  0 & -1\\
  1 & 1  
\end{pmatrix} \quad
B=\begin{pmatrix}
    1 & 1\\
    -1 & 0  
  \end{pmatrix}
\]
Consider $A$\\
$A^5=I$ thus $\ang{AI}=\ang{A}$ thus $\ang{A}$ is a subgroup of order 5\\
Consider $B$\\
$B^5=I$ thus $\ang{AI}=\ang{B}$ thus $\ang{B}$ is a subgroup of order 5\\
Since $A \neq B$ and $-A \neq B$ we know that these are distinct in $PSL_2(F_5)$\\
Thus since a group of order 60 has more than 1 subgroup of order 5 we know it must have 5 by Sylow 3, and thus it must be isomorphic to $A_5$

\newpage
\subsection*{Problem 4}
\begin{enumerate}[(a)]
    \item Consider $SL_2(F_7) = \{\t{2x2 matrices } | \t{ det} = 1\}$\\
    Thus for matrix $\begin{pmatrix}
        a & b\\
        c & d
    \end{pmatrix}$ in $SL_2(F_7)$ we know $ad-bc=1$\\
    Consider the first column we have $7^2-1$ choices since all it has to be is non zero\\
Then for the second column we have $7$ choices for $b$ but then $d$ is fixed to satisfy $ad-bc=1$ thus we have $120$ choices thus $|SL_2(F_7)|=7(7^2-1)=336$\\
Thus since $PSL_2(F_7)=\dfrac{SL_2(F_7)}{\{\pm I\}}$ we know $|PSL_2(F_7)|=\dfrac{|SL_2(F_7)|}{|\{\pm I\}|}=168$
\item Let $x$ and $y$ be two classes of matrices
\[
A=\begin{pmatrix}
  1 & 1\\
  0 & 1  
\end{pmatrix} \quad
B=\begin{pmatrix}
    1 & 0\\
    1 & 1  
  \end{pmatrix}
\]
Consider $A$\\
$A^7=I$ thus $\ang{AI}=\ang{A}$ thus $\ang{A}$ is a subgroup of order 7\\
Consider $B$\\
$B^7=I$ thus $\ang{AI}=\ang{B}$ thus $\ang{B}$ is a subgroup of order 7\\
Since $A \neq B$ and $-A \neq B$ we know that these are distinct in $PSL_2(F_7)$\\
Thus since a group of order 168 has more than 1 subgroup of order 7 we know it must have 8 by Sylow 3.
\item Let $N \unlhd G$ s.t. $[G : N]=2$ this means that $|N|=84$ thus $n_7=1$ by Sylow 3. However, this is a contradiction since proven in class if $P \unlhd N \unlhd G$ and $P$ is a $p$-Sylow then $P \unlhd G$ this implies that the one normal
seven Sylow is normal in $G$, but this is a contradiction of Sylow 2 thus it can't be index 2.
\item Suppose $[G : N]=3$ this means $|N|=56$ thus $n_7=1,8$ by Sylow 3. It must have 8 by shown above thus it has all 7-Sylow subgroups of $G$. Thus, it must contain both $H$ and $K$\\
Thus by group closure it contains element $u=xy^2x$ s.t. $H=\ang{x}$ and $K=\ang{y}$\\
$u=AB^2A$\\
$u^2=xy^2xxy^2x=xy^2x^2y^2x$\\
$u^3=xy^2x^2y^2xxy^2x=xy^2x^2y^2x^2y^2x=AB^2A^2B^2A^2B^2A=I$\\
This is a contradiction since $N$ is a group of order 56 and 3 does not divide 56 thus it cannot have a subgroup of order 3 thus $[G : N] \neq 3$
\item Assume $[G : N] \neq 6$ thus $|N|=28$ thus by Sylow 3 $n_7=1$ this is a contradiction by similar proof above thus $[G : N] \neq 6$
\item Assume $G$ is not simple and let $N \unlhd G$ be a maximal proper normal subgroup this means that $G \diagup N$ must be simple thus we only have to check for groups of prime order that divide 168 and since $60$ does not divide 168 we don't have to check this\\
Thus $[G:N]=7$ since we checked for 2,3 thus $|N|=24$
\item Thus by Sylow 3 we know that $n_2(N)=1,3$\\
Assume $n_2(N)=1$ then we know there is a normal 2-Sylow $H$ of $N$ that is normal in $G$. Thus, $G \diagup H$ is a subgroup of order 21 since $168/8=21$ thus by Sylow 3 this has a normal 7-Sylow $\bar{K}$. Thus, by proven in homework we know
corresponds to a normal subgroup in $G$ called $K$ of $G$ thus since $|K /diagup  H|=\bar{K}$ we know that $|K|=8\times 7=56$\\
Thus $[G : K]=3$ which is a contradiction thus $n_2(N)=3$\\
\item Thus $G$ acts on the 3 groups 2-Sylow groups by conjugation thus there exists a nontrivial homomorphism $\varphi : G \rightarrow S_3$. This map cannot be injective since  $|G|>|S_3|$ thus the
kernel is nontrivial, thus ker$(\varphi)=K$ is a proper normal subgroup of $G$.\\
By the first isomorphism theorem we know that $G \diagup K \cong \t{Im}(\varphi)=I$ and the $ |I| \mid 6$ and is not 1 thus we get $|I|=3,6$ thus $[G : K]=|I|=3,6$ and above we have shown these are contradictions thus $G$ must be simple
\end{enumerate}
\end{document}